%!TEX root = ../../exa-ma-d7.1.tex
\section{Software: FUnTiDES}
\label{sec:FUnTiDES:software}

\paragraph{Developer review.}
Aurélien Citrain and Hélène Barucq: please check the scientific scope,
supported configurations and software policies below, and add the reference
publications. Application and benchmark associations will be agreed separately.

\subsection{Purpose and scientific scope}
FUnTiDES (Fast Unstructured Time Dynamic Equation Solver) is a library for
two- and three-dimensional acoustic and elastic wave propagation, developed
by the Inria Makutu team. It combines a C++17 computational backend with
Python bindings and provides high-order spectral element methods (SEM)....

\subsection{Catalogue information}
\label{sec:FUnTiDES:metadata}
% Generated by gen.py; edit the catalogue snapshot, not this file.
{\small
\setlength{\LTleft}{0pt}
\setlength{\LTright}{0pt}
\renewcommand{\arraystretch}{1.12}
\begin{longtable}{@{}p{0.30\linewidth}p{\dimexpr0.70\linewidth-2\tabcolsep\relax}@{}}
\toprule
\textbf{Characteristic} & \textbf{Catalogue entry} \\
\midrule
\endhead
Framework & FUnTiDES \\
Exa-MA partner & Inria BXSO \\
Contacts & aurelien.citrain@inria.fr, helene.barucq@inria.fr \\
Licence category (survey) & OSS: Cecill-* \\
Architecture category (survey) & CPU or GPU \\
Bottlenecks & B7 - Exascale Algorithms \\
Languages & C++17, Python \\
Parallel programming & MPI, Multithread - OpenMP, GPU - Cuda, Kokkos \\
Data formats & Adios, Json, YAML, ASCII \\
Development practices & Continuous Integration, Container - Docker, Test - Unit, Test - Validation \\
Training resources & Documentation \\
Source repository & \href{https://github.com/FUnTiDES-sim/FUnTiDES}{Source repository} \\
Documentation & \href{https://github.com/FUnTiDES-sim/FUnTiDES\#readme}{Documentation} \\
Training material & \href{https://github.com/FUnTiDES-sim/FUnTiDES\#quick-start-build-and-run}{Training material} \\
Support & \href{https://github.com/FUnTiDES-sim/FUnTiDES/issues}{Support} \\
WP1 topics & finite element, spectral element, unstructured mesh, dG/hdG \\
WP7 topics & Geophysics, proxy-apps \\
\bottomrule
\end{longtable}
}
\noindent\textit{Source:} \href{https://docs.google.com/spreadsheets/d/19v57jpek52nQV2V0tBBON5ivGCz7Bqf3Gw-fHroVHkA/edit?gid=0}{Frameworks, row 14}; retrieved 2026-09-11.
The survey's architecture category does not establish that an Exa-MA
benchmark campaign has been completed.

\noindent No WP benchmark flag is checked in this snapshot.
Unchecked flags do not establish the absence of upstream tests or benchmarks.


\subsection{Implementation and available features}


\paragraph{Documentation and training.}


\subsection{Position in Exa-MA}
The catalogue identifies discretization topics for WP1 and geophysics
and proxy-app topics for WP7, with B7 (Exascale Algorithms) as the
reported bottleneck.
At this stage, no Exa-MA application, mini-app, extended mini-app,
demonstrator, proxy-app or benchmark record is associated with FUnTiDES.
The upstream project's use of the term \emph{proxy application} does not
automatically qualify an example as an Exa-MA proxy-app under the
project's application taxonomy.

\paragraph{Points to complete with the developers.}
Confirm the intended WP contributions (WP1, WP7?) and tested CPU/GPU/MPI configurations;
clarify consortium membership, resilience mechanisms, API stability and
versioning/metadata policies; supply the preferred publications or software
citation. These fields are not supplied in the catalogue snapshot, rather
than known to be absent from the software.

